\documentclass[12pt]{article}


\usepackage[a4paper, margin=1in]{geometry}
\usepackage{amsmath, amssymb}
\usepackage{bm}
\usepackage{setspace}
\usepackage{etoolbox}


\begin{document}


\onehalfspacing


\setlength{\parskip}{0.4em}
\setlength{\parindent}{0pt}
\setlength{\itemsep}{0.2em}


\AtBeginEnvironment{equation}{\vspace{0.2cm}}
\AtEndEnvironment{equation}{\vspace{0.2cm}}
\AtBeginEnvironment{align}{\vspace{0.2cm}}
\AtEndEnvironment{align}{\vspace{0.2cm}}


%------------------------------------------------


\begin{center}
\textbf{CSE 400: Fundamentals of Probability in Computing}\\
Lecture 20 Scribe Submission


\vspace{0.4cm}


Name: Aryan Tarachandani\\
Enrollment No: AU2420177\\
Email: aryan.t1@ahduni.edu.in\\
Date of Submission: March 26, 2026
\end{center}


\vspace{0.4cm}
\noindent\rule{\textwidth}{0.5pt}
\vspace{0.6cm}


%------------------------------------------------


\section*{Lecture 20: Linear Transformations and Expectations of Multiple RVs}


\vspace{0.3cm}


%------------------------------------------------


\section*{1. Random Vector Representation}


\vspace{0.2cm}


For random variables $X_1, X_2, \dots, X_N$, we represent them as:


\[
X =
\begin{bmatrix}
X_1 \\
X_2 \\
\vdots \\
X_N
\end{bmatrix}
\]


\begin{itemize}
\item $X$ is called a random vector
\item It provides a compact representation of multiple random variables
\end{itemize}


%------------------------------------------------


\section*{2. Mean Vector}


\vspace{0.2cm}


\textbf{Definition:}


\[
\mu_X = E[X]
\]


\begin{itemize}
\item Represents the average value of the random vector
\item Each element corresponds to the mean of one random variable
\end{itemize}


\textbf{Structure:}


\[
\mu_X =
\begin{bmatrix}
E[X_1] \\
E[X_2] \\
\vdots \\
E[X_N]
\end{bmatrix}
\]


%------------------------------------------------


\section*{3. Correlation Matrix}


\vspace{0.2cm}


\textbf{Definition:}


\[
R_{XX} = E[XX^T]
\]


\[
R_{XX}(i,j) = E[X_i X_j]
\]


\textbf{Matrix Form:}


\[
R_{XX} =
\begin{bmatrix}
E[X_1X_1] & E[X_1X_2] & \cdots & E[X_1X_N] \\
E[X_2X_1] & E[X_2X_2] & \cdots & E[X_2X_N] \\
\vdots & \vdots & \ddots & \vdots \\
E[X_NX_1] & E[X_NX_2] & \cdots & E[X_NX_N]
\end{bmatrix}
\]


\textbf{Property:}


\begin{itemize}
\item $R_{XX}$ is a symmetric matrix
\end{itemize}


%------------------------------------------------


\section*{4. Covariance Matrix}


\vspace{0.2cm}


\textbf{Definition:}


\[
C_{XX} = E[(X - \mu_X)(X - \mu_X)^T]
\]


\textbf{Matrix Form:}


\[
C_{XX} =
\begin{bmatrix}
Var(X_1) & Cov(X_1,X_2) & \cdots \\
Cov(X_2,X_1) & Var(X_2) & \cdots \\
\vdots & \vdots & \ddots
\end{bmatrix}
\]


\textbf{Property:}


\[
Cov(X_i, X_j) = Cov(X_j, X_i)
\]


%------------------------------------------------


\section*{5. Linear Transformation of Random Vectors}


\vspace{0.2cm}


\textbf{Model:}


\[
Y = AX + b
\]


\begin{itemize}
\item $A$: Transformation matrix (scales and rotates the data)
\item $b$: Constant vector (shifts/translates the center)
\item $X$: Input random vector
\end{itemize}


\textbf{Expanded Form:}


\[
Y_1 = a_{11}X_1 + a_{12}X_2 + \cdots + a_{1N}X_N + b_1
\]


\[
Y_2 = a_{21}X_1 + a_{22}X_2 + \cdots + a_{2N}X_N + b_2
\]


\[
\vdots
\]


\[
Y_N = a_{N1}X_1 + a_{N2}X_2 + \cdots + a_{NN}X_N + b_N
\]


%------------------------------------------------


\section*{6. Mean of Transformed Vector}


\vspace{0.2cm}


\textbf{Given:}


\[
Y = AX + b
\]


\textbf{Find:}


\[
\mu_Y = E[Y]
\]


\textbf{Derivation:}


\[
\mu_Y = E[AX + b]
\]


\[
= E[AX] + E[b]
\]


\[
= A E[X] + b
\]


\textbf{Final Result:}


\[
\mu_Y = A\mu_X + b
\]


%------------------------------------------------


\section*{7. Correlation Matrix Transformation}


\vspace{0.2cm}


\textbf{Definition:}


\[
R_{YY} = E[YY^T] = E[(AX + b)(AX + b)^T]
\]


\textbf{Expansion:}


\[
R_{YY} = E[AXX^TA^T + AXb^T + bX^TA^T + bb^T]
\]


\textbf{Applying Expectation:}


\[
R_{YY} = A E[XX^T] A^T + A E[X] b^T + b E[X^T] A^T + bb^T
\]


\textbf{Final Result:}


\[
R_{YY} = A R_{XX} A^T + A \mu_X b^T + b \mu_X^T A^T + bb^T
\]


%------------------------------------------------


\section*{8. Covariance Transformation}


\vspace{0.2cm}


\textbf{Key Observation:}


\[
Y - \mu_Y = (AX + b) - (A\mu_X + b)
\]


\[
= A(X - \mu_X)
\]


\begin{itemize}
\item The deviation from the mean is independent of the shift $b$
\end{itemize}


%------------------------------------------------


\section*{9. Case Study: Smart City Traffic + Weather System}


\vspace{0.2cm}


\textbf{Problem:}


Design an intelligent system in Ahmedabad that predicts traffic congestion using transportation and climate variables.


%-----------------


\textbf{Step 1: Random Vector}


\[
X =
\begin{bmatrix}
X_1 \\
X_2 \\
X_3 \\
X_4 \\
X_5
\end{bmatrix}
\]


\begin{itemize}
\item N-dimensional random vector with $N = 5$
\end{itemize}


%-----------------


\textbf{Step 2: Mean Vector}


\[
\mu_X =
\begin{bmatrix}
E[X_1] \\
E[X_2] \\
E[X_3] \\
E[X_4] \\
E[X_5]
\end{bmatrix}
\]


%-----------------


\textbf{Step 3: Covariance Interpretation}


\begin{itemize}
\item More rain $\Rightarrow$ more congestion
\item More rain $\Rightarrow$ lower speed
\item More traffic $\Rightarrow$ higher pollution
\item Higher temperature $\Rightarrow$ pollution disperses
\item Covariance matrix = dependency map of the city
\end{itemize}


%-----------------


\textbf{Step 4: Joint Gaussian Assumption}


\begin{itemize}
\item The random vector follows a joint Gaussian distribution
\item The system behaves in a joint probabilistic manner
\end{itemize}


%-----------------


\textbf{Step 5: Linear Transformation}


\[
Y = AX + b
\]


\begin{itemize}
\item Produces a congestion index
\end{itemize}


%-----------------


\textbf{Inference Statements}


\begin{itemize}
\item Rainfall increases congestion and decreases speed
\item Traffic increase leads to higher pollution
\item Temperature affects pollution dispersion
\end{itemize}


%-----------------


\textbf{Cause--Effect Relationships}


\begin{itemize}
\item Rain $\uparrow \Rightarrow$ Speed $\downarrow$
\item Rain $\uparrow \Rightarrow$ Traffic $\uparrow$
\item Traffic $\uparrow \Rightarrow$ Pollution $\uparrow$
\item Temperature $\uparrow \Rightarrow$ Pollution $\downarrow$
\end{itemize}


%------------------------------------------------


\end{document}